\section{Ergebnisse}
\label{sec:ergebnisse}

\subsection{Ausbeute}
Theoretische Ausbeute: $n_{\text{theo}} = \SI{0.10}{\mol}$.
Erhalten: \SI{6.2}{\gram} ($n = \SI{0.070}{\mol}$), also
\begin{equation}
  \text{Ausbeute} = \frac{n_{\text{erhalten}}}{n_{\text{theo}}} = \SI{70}{\percent}
  \label{eq:ausbeute}
\end{equation}

\subsection{Identifikation und Reinheit}
\begin{itemize}
  \item Siedepunkt: \SI{101}{\celsius} (Literatur: \SI{102}{\celsius}).
  \item DC: $R_f = 0.62$, keine Nebenprodukte sichtbar.
  \item IR (cm$^{-1}$): breite Bande 3350 ($\nu_{\mathrm{O-H}}$),
        2970 ($\nu_{\mathrm{C-H}}$), keine C=C-Bande bei 1640.
\end{itemize}

\subsection{NMR-Daten (Auszug)}
\begin{table}[H]
  \centering
  \caption{\textsuperscript{1}H-NMR (400 MHz, CDCl\textsubscript{3}).}
  \label{tab:nmr}
  \begin{tabular}{llll}
    \toprule
    $\delta$ in \si{ppm} & Multiplizität & $n$ H & Zuordnung \\
    \midrule
    1.20 & s & 6H & 2 $\times$ CH\textsubscript{3} \\
    1.45 & q & 2H & CH\textsubscript{2} \\
    0.90 & t & 3H & CH\textsubscript{3} (terminal) \\
    \bottomrule
  \end{tabular}
\end{table}
